\section{Introduction}

Multi-agent systems built on large language models (LLMs) are increasingly deployed across domains requiring autonomous task decomposition, long-horizon reasoning, and dynamic delegation. As these systems scale—spawning sub-agents, cascading tasks across service boundaries, and coordinating across distributed contexts—the problem of \textit{autonomous drift} emerges as a critical failure mode. Drift occurs when an agent or sub-agent, through successive autonomous spawning events, becomes semantically or operationally decoupled from its originating context, losing fidelity to the original intent while continuing to execute as if aligned. In production environments, drift manifests as任务漫游(task wander), where spawned agents pursue locally coherent but globally divergent objectives, or as \textit{authority escalation}, where sub-agents iteratively expand their own permissions in ways the root principal never sanctioned.

The severity of drift is compounded by the opacity of modern agentic pipelines. Unlike monolithic systems, multi-agent architectures built on frameworks such as AgentOS~\cite{bxthre3:agentos} decompose problems into hierarchical task graphs where each node is a potentially long-running autonomous process. When a root agent spawns a child agent to handle a subroutine, and that child further spawns tertiary agents, the resulting lineage is a branching execution tree whose leaves may be dozens of hops removed from the root context. Without structural enforcement mechanisms, the original constraints, preferences, and termination conditions encoded at the root attenuate at each generation, rendering downstream agents increasingly ungovernable.

\subsection{Why Existing Approaches Are Insufficient}

The canonical countermeasures against drift in existing multi-agent systems are time-to-live (TTL) constraints and depth-limited spawning. TTL-based approaches assign each agent a fixed computational or temporal budget; when the budget expires, the agent is terminated regardless of task state. Depth-limited approaches impose a maximum generation cap on the spawning tree, such that agents at depth $d_{\max}$ are prohibited from spawning further children. Both methods areappealing in their simplicity, but they are fundamentally insufficient as standalone drift controls.

TTL suffers from \textit{deadline arbitrariness}: there exists no principled method for assigning a TTL that is simultaneously short enough to bound drift risk and long enough to permit task completion. For open-ended reasoning tasks, TTL violations are an asymptotic failure mode—the agent approaches correct answers asymptotically but never arrives, while the TTL forces premature termination. Conversely, tasks that converge rapidly are penalized by TTLs sized for worst-case paths. Furthermore, TTL provides no semantic awareness; an agent approaching a correct conclusion may be killed while an agent diverging into a hallucinated subtask may be permitted to continue until its TTL expires. TTL is a blunt instrument that controls \textit{duration} without controlling \textit{direction}.

Depth-limited spawning addresses the attenuation problem by hard-stopping lineage growth, but it introduces its own failure mode: \textit{depth coercion}. Many tasks require arbitrarily deep decomposition—formal verification, multi-step code generation, recursive theorem proving—in ways that are semantically necessary and not reducible by clever prompt engineering. A depth limit of $d=5$ is arbitrary relative to a task that requires $d=6$, and there exists no context-independent method for calibrating $d_{\max}$ across heterogeneous task distributions. Moreover, depth limits do not prevent drift \textit{within} a generation: an agent at depth $d$ may still spawn multiple children that collectively drift from the root intent while remaining within the depth budget.

A third class of approaches, constraint-passing schemes, attempt to propagate root-level invariants (preconditions, postconditions, budget limits, ethical boundaries) through the spawning chain. However, constraint-passing is fragile by construction: agents may reinterpret, underspecify, or outright discard constraints they find inconvenient or incomprehensible. Without structural enforcement, constraint-passing relies on the good faith of autonomous agents—an unreliable assumption for systems operating in adversarial or high-stakes environments.

\subsection{Core Insight: Immutable Parent Pointer Chain}

The fundamental limitation of TTL, depth limits, and constraint-passing schemes is that they address drift \textit{after} it has already occurred, through outcome-based termination or attenuation, rather than \textit{preventing} the structural conditions that enable drift. The core insight of the Recursive Spawning framework is that drift is structurally impossible in a spawning tree where every agent maintains an \textit{immutable reference pointer} to its parent, and where spawning is defined as a referentially transparent operation with respect to that pointer chain.

Formally, let each agent $a_i$ in a spawning tree be associated with an immutable parent pointer $p_i$ such that $p_i \mapsto a_{\text{parent}(i)}$, where $\text{parent}(i)$ denotes the generating agent in the preceding generation. Because $p_i$ is immutable after creation, an agent cannot be ``reparented'' to a new context, cannot sever its connection to the root, and cannot create a child whose parent pointer chain does not ultimately trace back to the root principal. Every node in the tree is thus a \textit{provable descendant} of the originating context, and every operation performed by any node remains referentially anchored to that context.

This structural invariant yields a powerful consequence: an agent cannot drift in any way that is not detectable, attributable, and attributable to a specific point in the parent pointer chain. Drift, by definition, requires semantic or operational decoupling from the originating context. But if the parent pointer is immutable and the chain is complete, decoupling is structurally impossible—there is no mechanism by which an agent can ``forget'' its provenance. The agent may still deviate in its \textit{contentual reasoning}, but it cannot deviate in its \textit{structural position} within the lineage graph. Drift that is definitionally tied to structural decoupling becomes, under immutable parent pointers, a logical impossibility.

This reframes the drift problem from one of \textit{agent behavior control}—hard to solve, brittle, and requiring constant monitoring—to one of \textit{tree structure enforcement}, which is achievable through immutable references and referential transparency guarantees. Recursive Spawning thus replaces behavioral policing with structural invariance: instead of asking ``will this agent obey its constraints?'', the framework asks ``is this agent provably connected to its root?''. The latter question has a verifiable, stable answer; the former does not.

\subsection{Formal Definition of Drift}

To make the above precise, we define drift formally within the BX3 agentic framework. Let $\mathcal{A}$ denote the set of all agents in a spawning tree, and let $\rho: \mathcal{A} \to \mathcal{A}^*$ map each agent to the ordered sequence of its ancestors, from immediate parent to root. That is, for agent $a$, $\rho(a) = (a_{\text{parent}}, a_{\text{grandparent}}, \ldots, a_{\text{root}})$.

\textbf{Definition 1 (Drift).} An agent $a \in \mathcal{A}$ is said to have \textit{drifted} if there exists some predicate $\phi$ encoding the root intent such that $\phi(a)$ is false while for all $a' \in \rho(a)$ (the ancestor chain), $\phi(a')$ is true, and this violation is not attributable to a correct semantic decision by $a$ but rather to a loss of fidelity to the originating context.

This definition distinguishes drift from legitimate divergence: an agent may correctly reason that a sub-task requires a strategy different from what its ancestors anticipated. Such divergence is not drift because it is a \textit{semantic correction}, not a fidelity loss. Drift is specifically the case where the agent's behavior violates a constraint or objective that was present and operative at the root, without this violation being justified by superior local information.

The Recursive Spawning framework guarantees that structural drift (as distinct from semantic divergence) cannot occur, because the immutable parent pointer chain ensures that every agent remains referentially anchored to the root. Semantic drift—where an agent makes a locally wrong decision—remains possible, but it is within the scope of existing HITL (human-in-the-loop) and truth-gating mechanisms in the BX3 stack, which are designed precisely to catch and correct semantic errors at task completion boundaries.

\subsection{Paper Roadmap}

The remainder of this paper is organized as follows. Section~\ref{sec:framework} presents the Recursive Spawning framework in detail, including the immutable parent pointer protocol, the spawn operation semantics, and the structural invariants that enforce drift impossibility. Section~\ref{sec:implementation} describes the integration of Recursive Spawning within the BX3 AgentOS stack, covering the Rust/TypeScript backend, the event cascade backstop, and the SSE streaming interfaces used for real-time lineage monitoring. Section~\ref{sec:experiments} presents empirical evaluation of the framework on benchmark task suites measuring drift frequency, task completion rates, and cascade reliability against baseline TTL and depth-limited approaches. Section~\ref{sec:related} surveys related work in multi-agent coordination, autonomous governance, and LLM-based task decomposition. Section~\ref{sec:conclusion} concludes with a discussion of limitations, open questions, and the path toward production-scale deployment of Recursive Spawning across the BX3 portfolio.
